% !PRINTED 512
\weberpage{512}

Daraus schliessen wir noch nach der Formel (1), dass \gesperrt{associirte Functionale dieselbe absolute Norm haben}.

Ein ganzes rationales Functional ist hiernach immer mit seinem absoluten Werthe associirt. Ist also $\alpha$ irgend ein ganzes Functional des Körpers $\overline{\Omega}$, so sind auch $N(\alpha)$ und $N_a(\alpha)$ associirt. Da nun in $N(\alpha)=\alpha\alpha'$ der Factor $\alpha'$ als Product von ganzen Functionalen selbst ganz ist, so ist $N(\alpha)$ und folglich auch die natürliche Zahl $N_a(\alpha)$ durch $\alpha$ theilbar. Wir formuliren also noch den Satz:

13. \gesperrt{Es giebt natürliche ganze Zahlen (in unendlicher Menge), darunter die absolute Norm von $\alpha$, die durch ein beliebiges ganzes Functional $\alpha$ theilbar sind. Ist $a$ die kleinste unter den durch $\alpha$ theilbaren natürlichen Zahlen, so ist jede durch $\alpha$ theilbare ganze rationale Zahl durch $a$ theilbar.}

Denn ist $m$ eine durch $\alpha$ theilbare ganze rationale Zahl, so ist auch der Rest der Division von $m$ durch $a$ eine durch $\alpha$ theilbare, ganze rationale Zahl. Da diese kleiner als $a$ ist, so muss sie $=0$ sein.

\S.\,139.

\gesperrt{Grösster gemeinschaftlicher Theiler.}

Es mögen $\alpha, \beta, \ldots$ von Null verschiedene ganze Functionale in $\Omega$ in beliebiger aber endlicher Anzahl bedeuten, und $x, y, \ldots$ Variable, die in $\alpha, \beta, \ldots$ nicht vorkommen. Dann ist

\weq{1}{\delta = \alpha x + \beta y + \cdots}

gleichfalls ein ganzes Functional in $\Omega$. Die Norm von $\delta$ ist eine ganze homogene Function $n^{\text{ten}}$ Grades der Variablen $x, y, \ldots$, deren Coëfficienten ganze rationale Functionale sind. Bezeichnen wir die absolute Norm von $\delta$ mit $D$, so ist

\weq{2}{N(\delta) = D\,E(x, y, \ldots) = D\,E,}

und darin ist $E(x, y, \ldots) = E$ eine ganze rationale Function der Variablen $x, y, \ldots$ und im Allgemeinen eine gebrochene Function der in $\alpha, \beta, \ldots$ vorkommenden Variablen, jedenfalls aber eine \gesperrt{rationale Einheit} [\S.\,136, (2)].
